% !TEX root = ../dissertationOG.tex

\clearpage % clear the prior chapter's page

\chapter{Related Work}\label{CH2_RelatedWork}
%\vspace{-7mm}
%\bigskip

\section{Physics-Based Shape Morphing}
% ============================================================================
\subsection{Material Point Method (MPM)}
% ============================================================================
    The Material Point Method (MPM) was introduced by~\citep{sulsky1994particle} to solve solid mechanics problems involving large deformations and fracture. As a hybrid method that couples Lagrangian particles with an Eulerian grid, MPM effectively handles severe material deformation and topology-changing phenomena.
    
    Early MPM implementations that relied on standard PIC/FLIP transfers suffered from noticeable numerical dissipation and, in particular, poor conservation of angular momentum. To mitigate this, the \textbf{APIC} (Affine Particle-in-Cell) scheme~\citep{jiang2017angular} was proposed; by transporting an affine velocity term ($\mathbf{C}$), APIC better preserves angular momentum and shear response, thereby reducing artificial damping. Building on this idea, \textbf{MLS-MPM}~\citep{2018-MLSMPM} leverages Moving Least Squares to robustly estimate the deformation gradient ($\mathbf{F}$), alleviating the classic cell-crossing instabilities when particles traverse grid boundaries and improving stability under extreme deformations.
    
    Recent work has substantially extended MPM in stability, accuracy, and scope. Implicit time-integration techniques~\citep{implicit_mpm_surface_tension, wang2020hierarchical} increase robustness for stiff materials and allow larger time steps, while high-fidelity frictional contact algorithms~\citep{kakouris2024high} are now available for scenarios with repeated contact and separation, such as the cutting task considered in this thesis.
    
    The repertoire of material models has likewise broadened from simple elasticity to anisotropic media (e.g., meats and various fruits)~\citep{10.1145/3386569.3392428}, complex plasticity (e.g., sand, snow)~\citep{2019-cdmpm}, and fracture coupled with topological change~\citep{wang2020massively}. MPM has also been extended to challenging geometries, including thin-shell dynamics~\citep{guo2018material}.
    
    In parallel, performance-oriented research has accelerated MPM substantially: GPU acceleration with automatic differentiation (e.g., JAX-based pipelines)~\citep{du2025jax} and sparse, memory-efficient GPU-MPM frameworks for large-scale simulations~\citep{chen2025sparse} enable much larger problem sizes.
    
    \emph{Throughout this thesis, we build on these developments. In particular, Chapter~\ref{CH3_SHAPEMORPHING} will adopt APIC’s momentum-preserving transfers together with the stability of MLS-MPM to form a reliable computational core, which later chapters leverage for differentiable control.}

% ============================================================================
\subsection{Shape Morphing}
% ============================================================================
    Shape morphing algorithms have historically been explored under a variety of geometric representations. Implicit-surface approaches were demonstrated by~\citep{turk2005shape, breen2001level}; \citep{bansal2018lie} proposed interpolation via Lie body representations of triangle meshes; and~\citep{kilian2007geometric} re-cast shapes as points in a Riemannian shape space, enabling natural interpolation operations.
    
    Physically motivated approaches, which are closest in spirit to our goals, have also gained traction. For example,~\citep{bao2005physically} used point samples on surface meshes to obtain results with greater physical plausibility than linear blend methods, while~\citep{muller2004point} leveraged point samples for both surfaces and volumes to morph elastoplastic objects. Although Chapter~\ref{CH3_SHAPEMORPHING} shares this high-level objective, it \emph{critically differs} by employing an MPM-based physics simulation and by leveraging differentiable simulation and modern optimization capabilities that were not present in~\citep{muller2004point}.
    
    The widely used As-Rigid-As-Possible (ARAP) energy~\citep{2000-ARAPShapeInterpolation} and its derivatives deform triangle/tetrahedral meshes via energy minimization for interpolation tasks. Subsequent extensions have emphasized physical realism~\citep{2004-ActualMorphing, 2020-HamiltonianShapeInterpolation} to improve accuracy, e.g., solving a Poisson equation on the domain mesh~\citep{2005-PoissonShapeInterpolation} to reduce artifacts such as shrinkage. Further,~\citep{liu2011rigid} reduced reliance on consistent tetrahedral meshes by operating on triangle surface meshes, broadening practical applicability.
    
    More recently, data-driven methods have converged with geometry processing: \citep{2013-DataDrivenShapeMorphing} constructed local shape spaces from model databases and posed deformation as global optimization, while neural-field representations have enabled new routes to geometric editing~\citep{yang2021geometry}. Given the growing interplay between learning and simulation (see, e.g.,~\citep{sharpdeep}), \emph{Chapter~\ref{CH3_SHAPEMORPHING} will present a physics-based methodology that is designed to integrate with learning components in future work.}

% ============================================================================
\subsection{Optimization and Control}
% ============================================================================
Shape morphing is closely related to spatiotemporal optimization~\citep{1988-SpacetimeConstraints}, which directly informs our approach. Within this domain, prior work has investigated inverse control of physical systems:
\begin{itemize}
  \item \textbf{Fluid control:} McNamara et al.~\citep{2004-FluidControlAdjointMethod} minimized a keyframe-based objective by computing gradients via the adjoint method, aligning simulated fluid motion with target states.
  \item \textbf{Deformable-body control:} Coros et al.~\citep{2012-DeformableObjectsAlive} optimized nonlinear, time-dependent objectives with SQP, addressing issues such as rest-state adaptation during continuous control of mesh-based characters.
  \item \textbf{Amorphous-material manipulation:} Zhang et al.~\citep{2020-ManipulateAmorphousMaterialsRL} learned controllers for amorphous materials via reinforcement learning with externally controlled rigid contacts.
\end{itemize}
\emph{Chapter~\ref{CH3_SHAPEMORPHING} builds on this trajectory by framing shape morphing as an inverse control problem. Using differentiable MPM, we directly optimize physically meaningful parameters (e.g., deformation gradients) to reach a prescribed target state, in line with established spatiotemporal optimization principles.}

% ============================================================================
\subsection{Differentiable Simulation}
% ============================================================================
    \emph{The approach in Chapter~\ref{CH3_SHAPEMORPHING} is built on a differentiable simulation framework.} Differentiable simulation~\citep[see, e.g.,][]{coros2021differentiable} has attracted significant interest in the learning and graphics communities due to its flexibility and efficiency. Its core is the ability to compute \textbf{gradients} across the entire simulation via \textbf{backpropagation}, inspired by neural networks~\citep{2020-DiffPhysicsSim}.
    
    This paradigm has been applied to soft-robot control~\citep{2018-ChainQueen, 2019-DiffTaichi}, material parameter estimation~\citep{2019-DiffCloth, 2022-DiffSoftMultiBody}, and accelerated reinforcement learning~\citep{2021-PODS, 2022-DiffSimRL}. Among various discretizations, the \textbf{Material Point Method (MPM)}~\citep{sulsky1994particle, jiang2016material} has seen successful differentiable formulations and implementations (e.g.,~\citep{2018-ChainQueen}), which we leverage as a core component.
    
    Recent work has also addressed \textbf{collision handling} in complex dynamical systems within differentiable frameworks~\citep{qiao2020scalable}. Our methodology likewise maintains physically consistent constraints such as external forces. For instance, our morphing results remain elastically plausible under gravitational disturbances, underscoring the necessity and advantage of an MPM-based approach (see Section~\ref{sec:examples}).
    
    \emph{Chapter~\ref{CH3_SHAPEMORPHING} will further extend the applicability of differentiable MPM (including MLS-MPM~\citep{2018-MLSMPM}) to shape morphing with complex, \textbf{topology-aware} configurations---without the intricate mesh setups required by~\citep{10.1145/3130800.3130820}. Ultimately, this combination of simulation and animation techniques enables both smooth deformations and dynamic topology changes, as demonstrated in our results.}

\section{Gaussian Shape Morphing with a Differentiable Bridge}
% ============================================================================
\subsection{Scene Rendering with Gaussian Splatting}
% ============================================================================
    3D Gaussian Splatting (3DGS)~\citep{kerbl20233d} has emerged as the de facto standard for real-time, high-quality novel view synthesis. Scenes are represented as sets of anisotropic 3D Gaussians and rendered by a visibility-aware, differentiable rasterizer with explicit density control. However, in the traditional pipeline (simulation $\rightarrow$ \emph{surface extraction} $\rightarrow$ rendering), the threshold-based surface stage is \emph{non-differentiable}, which prevents pixel-space errors from flowing back to the physical state, leaving a persistent \emph{rendering gap}. \emph{Chapter~\ref{CH4_PhysMorphGS}} addresses this gap by \textbf{bypassing surface extraction} and enabling pixel losses to reach physical control variables directly.
    
    The success of 3DGS has spurred numerous extensions to dynamic scenes. Methods such as 4DGS~\citep{wu20244dgaussiansplattingrealtime} and Deformable 3D Gaussians~\citep{yang2023deformable3dgaussianshighfidelity} learn complex deformation fields or per-Gaussian trajectories to capture non-rigid motion, while \textbf{SplaTAM}~\citep{keetha2024splatam} targets SLAM for online, real-time reconstruction. Yet these motions are fundamentally \emph{kinematic}: they represent observed motion without explicit first-principles dynamics (forces, mass, material laws). \emph{Chapter~\ref{CH4_PhysMorphGS}} introduces a \textbf{particle$\rightarrow$Gaussian bridge} so that pixel losses backpropagate into the physical state, improving not only visual sharpness and temporal stability but also physical conservation metrics (mass, momentum, energy drift).
    
    Within this broader effort, PhysGaussian~\citep{xie2024physgaussian} achieves high fidelity by driving 3DGS animations with forward MPM simulation, but it remains \emph{unidirectional} (physics $\rightarrow$ render), blocking gradients from returning to physics. Conversely, OmniPhysGS~\citep{lin2025omniphysgs} replaces first-principles dynamics with \emph{neural surrogates}, gaining efficiency at the potential cost of physical-law preservation and generalization to unseen materials. \emph{Chapter~\ref{CH4_PhysMorphGS}} aims to inherit the strengths of both directions while closing the \emph{rendering gap} via a differentiable, surface-free coupling that completes a closed-loop optimization from pixels to physics.

% ============================================================================
\subsection{Differentiable Physics Coupled with Rendering}
% ============================================================================
    This resarch aims to unify \emph{first-principles differentiable physics} with \emph{photorealistic 3DGS rendering} in a single optimization loop. Early attempts using meshes or implicit representations~\citep{Yifan_2019,laine2020modular,wu2025pbr} struggled respectively with topology changes or computational complexity. PhysDreamer~\citep{zhang2024physdreamerphysicsbasedinteraction3d} also combines differentiable MPM with Gaussian rendering to fit pixel observations, but focuses on distilling video priors to learn a neural material field and initial velocities.
    
    In contrast, this research avoids learned priors and introduces a \textbf{fully differentiable particle$\rightarrow$Gaussian upsampling bridge} that allows pixel-space losses to directly update per-particle deformation states. Concretely, sparse MPM particle states $(x, F)$ are mapped to dense Gaussian attributes $(\mu, \Sigma)$ without explicit surface extraction. Silhouette, edge, and covariance losses ($L_\alpha$, $L_{\text{edge}}$, $L_{\text{cov}}$) then backpropagate via an adjoint path into the MPM state, establishing a \textbf{bidirectional coupling}. This design overcomes the limitations of forward-only pipelines, surrogate-only dynamics, or proxy-loss-only optimization and, critically, \textbf{reduces the rendering gap} by ensuring that pixel-level errors flow back to physical control variables.
    
{Chapter \ref{CH4_PhysMorphGS} (Rendering Gap).}
    This work systematically reduces the \emph{rendering gap} between differentiable simulation and 3DGS by replacing the non-differentiable surface extraction stage with a differentiable particle$\rightarrow$Gaussian bridge. Pixel losses directly reach the physical state (e.g., per-particle deformation gradients), and the loop is designed to improve visual metrics (boundary sharpness, temporal stability) alongside physical conservation metrics (mass, linear/angular momentum, energy drift), aligning visual coherence with physical fidelity in a closed-loop manner.

\section{Force-Aware Cutting (Verifying Physical Interaction with MPM)}
% ============================================================================
\subsection{Large Multimodal Models (LMMs).}
% ============================================================================
    Text-only large language models (LLMs) exhibit strong reasoning and generalization, motivating multimodal extensions that accept images and video~\citep{achiam2023gpt}. A common recipe aligns a frozen visual encoder to an open LLM via a lightweight projector or adapter, enabling instruction-following in vision--language dialogue~\citep{radford2021clip,liu2023llava,zhu2023minigpt4}. Beyond holistic image understanding, region-level LMMs improve referential precision by accepting spatial prompts or RoI features, supporting pointing, boxes, and segmentation during conversation~\citep{cai2023vipllava,rasheed2024glamm,chen2023shikra,zhang2024gpt4roi}. Meanwhile, video LMMs extend temporal reasoning with interleaved visual--text tokens and long-context handling~\citep{ataallah2024minigpt4video,li2024llavainterleave}. To systematically assess these capabilities, recent benchmarks probe general multimodal comprehension (MMBench), hierarchical and text-rich settings (SEED-Bench / SEED-Bench-2), and long-form temporal understanding (MMBench-Video)~\citep{liu2024mmbench,li2023seedbench,li2023seedbench2,fang2024mmbenchvideo}. Finally, MAGMA-style adapter methods foreshadow lightweight multimodal augmentation~\citep{eichenberg2022magma}, while the more recent Magma agent model introduces \emph{Set-of-Mark} and \emph{Trace-of-Mark} as surrogate tasks to bridge verbal understanding with spatio--temporal action across UI navigation and robotics~\citep{yang2025magma}. \emph{In this context, Chapter \ref{CH5_CUTTINGSIM} goes beyond purely visual--linguistic alignment by incorporating \textbf{force-aware} signals (impulse/force traces) into the learning loop.}

% ============================================================================
\subsection{Vision--Language--Action for Robotics.}
% ============================================================================
    Several studies adapt large multimodal models to robot control~\citep{rt2,openvla,prismaticvlm,llarva,tracevla,lapa,moto}. RT-2 shows that co-finetuning on web-scale vision--language tasks and robot trajectories, while decoding discretized action tokens, can ground natural-language instructions into low-level actions~\citep{rt2}. OpenVLA standardizes an open VLA stack by building on an open-source Prismatic VLM backbone and releasing a reusable action interface, enabling stronger generalization under a public, reproducible recipe~\citep{openvla,prismaticvlm}. To tighten spatial grounding and planning, LLARVA predicts 2D visual traces as an auxiliary target to align vision and action spaces~\citep{llarva}, while TraceVLA uses visual trace prompting to inject spatio--temporal intent into action generation and improves generalist policies across embodiments~\citep{tracevla}. Learning directly from videos via discrete latent tokenization has also emerged as a scalable pretraining signal for control, either by predicting VQ-VAE--based latent actions or motion tokens before mapping to robot actuation~\citep{lapa,moto}. \emph{By contrast, Chapter \ref{CH5_CUTTINGSIM} adopts a \textbf{physics-consistent VLA approach} that preserves OpenVLA-style action interfaces while coupling them with \textbf{SoM/ToM supervision} and an \textbf{MPM-based force feedback} loop.}
    
% ============================================================================
\subsection{Cutting-focused framing.}
% ============================================================================
    Cutting requires precise initiation on object surfaces, stable contact through style phases such as sawing and press-through, and compliance to material properties. For compatibility with existing VLA infrastructures, we follow OpenVLA’s action representation~\citep{openvla} and pair it with richer multimodal supervision describing cut locations, styles, and material cues. Instead of only predicting latent tokens from video~\citep{lapa,moto}, we supervise \emph{Set-of-Mark (SoM)} for \emph{where} to cut and \emph{Trace-of-Mark (ToM)} for \emph{how} the tool should move over time, which strengthens spatio--temporal reasoning for cutting and improves initiation accuracy, contact stability across style transitions, and cross-object generalization relative to prior VLA formulations~\citep{rt2,openvla,llarva,tracevla}. \emph{Chapter \ref{CH5_CUTTINGSIM} further aligns these intentions with \textbf{contact dynamics} by introducing a \textbf{force--time fidelity} objective (e.g., sDTW/peak/impulse terms) that ties SoM/ToM guidance to measured interaction forces.}
    
% ============================================================================
\subsection{Large-scale robot learning datasets.}
% ============================================================================
    Community efforts have consolidated broad, open datasets for imitation and VLA-style policies, but almost none target deformable cutting. \emph{Open X-Embodiment (RT-X)} aggregates over one million real-robot trajectories spanning 22 embodiments by pooling 60 prior datasets from 34 labs into a standardized RLDS/\texttt{tfrecord} format, and additionally releases pre-trained RT-X checkpoints to study cross-embodiment transfer~\citep{o2024open}. Beyond X-Embodiment, \emph{RoboNet} pioneered sharing multi-robot experience (about 15M video frames across seven platforms) for vision-based manipulation~\citep{robonet2019}. \emph{BridgeData~V2} contributes 60{,}096 trajectories across 24 environments on a low-cost platform, enabling goal/language conditioning and cross-institution generalization~\citep{bridgedata2023}. \emph{DROID} scales in-the-wild teleoperation to 76k demonstrations ($\sim$350 hours) over 564 scenes and 80+ tasks~\citep{khazatsky2025droidlargescaleinthewildrobot}. For evaluation, \emph{LIBERO} provides language-conditioned manipulation suites, but does not include deformable cutting~\citep{libero2023}. Cutting-specific resources are largely simulation-centric: \emph{DiSECt} offers a differentiable cutting simulator with small real datasets for calibration~\citep{disect2022}; \emph{SliceIt!} proposes a dual high-fidelity cutting+robot simulator with limited real slicing for real2sim2real transfer~\citep{sliceit2024}; and the recent \emph{TopoCut} introduces a benchmark for multi-step cutting with spectral rewards and diffusion policies~\citep{topocut2025}. \emph{These trends highlight a gap for large, multi-embodiment \textbf{real-robot} cutting datasets comparable in scale and standardization to X-Embodiment; Chapter \ref{CH5_CUTTINGSIM} addresses this by proposing a collection/evaluation protocol that includes ratio/direction conditioning, SoM/ToM annotations, F/T and pose labels, and a force-aware sim$\leftrightarrow$real loop built on MLS-MPM/CPIC.}
    